\problemname{Bitrussia}
\illustration{0.3}{coins.jpg}{Picture by \href{https://commons.wikimedia.org/wiki/File:Assorted_United_States_coins.jpg}{Elembis} (public domain)}
In the Republic of Bitrussia, a new coin system has recently been introduced.
There are $N$ different coin denominations worth $2^0, 2^1, 2^2, ..., 2^{N-1}$.

The small town of Napsaks is known for being full of interesting shops.
At the same time, Napsaks is notorious for never having any change in the shops.
You are also not allowed to pay more than the price.
It is therefore very important to bring plenty of coins of suitable denominations to be able to buy everything you want.

In Napsaks lives Darja-Pavla.
She is planning to go Christmas shopping and has brought $a_i$ coins of value $2^i$ ($i = 0, 1, ..., N-1$).
She will visit $M$ different shops and buy one thing in each shop.
The item she buys in shop $i$ costs $b_i$ ($i = 0, 1, ..., M-1$).
She is of course worried that her coins will not be enough to pay for everything she wants to buy. Help her determine this!

\section*{Input}
The first line contains two integers $1 \le N \le 50$ and $1 \le M \le 100\,000$, separated by spaces.
The next line contains the $N$ space-separated integers $0 \le a_0, a_1, ..., a_{N-1} \le 10^{15}$.
The third and last line contains the $M$ space-separated integers $0 \le b_0, b_1, ..., b_{M-1} \le 10^{15}$.

\section*{Output}
Output \texttt{ja} if Darja-Pavla can pay for everything she wants to buy with her coins.
Otherwise, output \texttt{nej}.

\section*{Scoring}
Your solution will be tested on a set of test groups. To earn points for a group, you must pass all test cases in that group.

\noindent
\begin{tabular}{| l | l | l |}
\hline
Group & Points & Constraints \\ \hline
1     & 19         & $M = 1, a_i \le 1$ \\ \hline
2     & 46         & $M = 1$ \\ \hline
3     & 35         & No additional constraints. \\ \hline
\end{tabular}

\section*{Explanation of Sample 1}
In the sample, Darja-Pavla has one coin worth $1$, three coins worth $2$, and one coin worth $4$.
She can pay for the item costing $5$ by using one coin worth $1$ and one coin worth $4$ ($1 + 4 = 5$).
She can then pay for the item costing $6$ with the three remaining coins of value $2$ ($2 + 2 + 2 = 6$).

\section*{Explanation of Sample 2}
In the sample, both items require that a coin of denomination $1$ is used, but she only has one coin of this denomination.
